
\documentclass{article}
\usepackage[dvipsnames]{xcolor}
\usepackage{graphicx} % Required for inserting images
\usepackage{geometry}
\usepackage{longtable}
\usepackage{amsmath}
\usepackage{amssymb}
\usepackage{booktabs}
% \usepackage{algorithm}
% \usepackage{algorithmic}

% \usepackage{algorithm}
\usepackage{algpseudocode}
\usepackage{booktabs}
\usepackage{siunitx}
\usepackage{graphicx}
\newgeometry{left=2cm,bottom=2cm}
\usepackage{amsthm}
\usepackage{xcolor}

\newtheorem{proposition}{Proposition}
\newtheorem{definition}{Definition}
\newtheorem{example}{Example}
% Define theorem environments
\newtheorem{theorem}{Theorem}
\newtheorem{corollary}{Corollary}
\newtheorem{lemma}{Lemma}

\title{Riemannian VAE}
\date{}

\begin{document}

\maketitle


{\color{red}
\textbf{Latent space:} The encoder produces $z \in B_T \subset \mathbb{R}^L$.

\textbf{Tangent space:} The decoder map $f_\theta$ takes $z$ and projects it into 
the tangent space at $\mu$:
\[
    f_\theta : B_T \to T_\mu M.
\]

\textbf{Manifold projection via exponential map:} The exponential map takes 
tangent vectors and maps it to points on the manifold $M$:
\[
    f^M_\theta = \mathrm{Exp}_\mu \circ f_\theta : B_T \to M.
\]
We can write $f_\theta^M(z)=\operatorname{Exp}_\mu\left(f_\theta(z)\right)$.
}




\section{Riemannian Geometry}
\subsection{(Sub)manifolds, Weighted (Sub)manifolds}

\begin{definition}[Riemannian manifold]
A Riemannian manifold \(M\) is a differentiable manifold \(M\) equipped with a Riemannian metric \(g\), which is a smoothly varying 
inner product on the tangent spaces of \(M\).
\end{definition}

\begin{example}[The 2-sphere]
Consider the 2-sphere
\[
S^2 = \{ (x,y,z) \in \mathbb{R}^3 : x^2 + y^2 + z^2 = 1 \}.
\]

\textbf{Step 1: Differentiable manifold.}  
\(S^2\) is a smooth 2-dimensional surface sitting in \(\mathbb{R}^3\).   Locally, we can describe it with coordinates, for instance spherical coordinates:
\[
(\theta,\phi) \mapsto (\sin\theta \cos\phi, \, \sin\theta \sin\phi, \, \cos\theta),
\]
with \(\theta \in (0,\pi)\) and \(\phi \in (0,2\pi)\).

\textbf{Step 2: Tangent spaces.}  
At each point \(p \in S^2\), the tangent space \(T_p S^2\) is the 2-dimensional plane in \(\mathbb{R}^3\) orthogonal to the radius vector \(p\).  

We define the Riemannian metric \(g\) as the restriction of the standard Euclidean inner product on \(\mathbb{R}^3\) to these tangent planes. Concretely, if 
\(\vec{u}, \vec{v} \in T_p S^2 \subset \mathbb{R}^3\), then
\[
g_p(\vec{u}, \vec{v}) = \vec{u} \cdot \vec{v}.
\]

\textbf{Step 4: Explicit form.}  
In spherical coordinates, this metric has the local expression
\[
g = d\theta^2 + \sin^2 \theta \, d\phi^2,
\]
which makes explicit how lengths and angles are measured on the surface of the sphere.


We parametrize the unit sphere by spherical coordinates:
\[
\vec{X}(\theta,\phi) = 
\bigl(\sin\theta \cos\phi,\ \sin\theta \sin\phi,\ \cos\theta\bigr),
\qquad \theta \in (0,\pi),\ \phi \in (0,2\pi).
\]

\textbf{Tangent basis:}  
The tangent basis vectors are
\[
\vec{e}_\theta = \frac{\partial \vec{X}}{\partial \theta} 
= (\cos\theta \cos\phi,\ \cos\theta \sin\phi,\ -\sin\theta),
\]
\[
\vec{e}_\phi = \frac{\partial \vec{X}}{\partial \phi} 
= (-\sin\theta \sin\phi,\ \sin\theta \cos\phi,\ 0).
\]

\textbf{Metric coefficients:}  
The induced metric is obtained from the Euclidean dot product:
\[
g_{ij} = \vec{e}_i \cdot \vec{e}_j.
\]
Hence
\[
g_{\theta\theta} = \vec{e}_\theta \cdot \vec{e}_\theta 
= \cos^2\theta(\cos^2\phi+\sin^2\phi) + \sin^2\theta = 1,
\]
\[
g_{\theta\phi} = \vec{e}_\theta \cdot \vec{e}_\phi 
= \cos\theta\cos\phi(-\sin\theta\sin\phi) 
+ \cos\theta\sin\phi(\sin\theta\cos\phi) + (-\sin\theta)\cdot 0 = 0,
\]
\[
g_{\phi\phi} = \vec{e}_\phi \cdot \vec{e}_\phi 
= \sin^2\theta(\sin^2\phi+\cos^2\phi) = \sin^2\theta.
\]

\textbf{Line element:}  
The infinitesimal displacement on the surface is
\[
d\vec{X} = \vec{e}_\theta\, d\theta + \vec{e}_\phi\, d\phi.
\]
Its squared length, using the Euclidean dot product, is
\[
ds^2 = d\vec{X}\cdot d\vec{X}
= (\vec{e}_\theta d\theta + \vec{e}_\phi d\phi)\cdot(\vec{e}_\theta d\theta + \vec{e}_\phi d\phi)
= \begin{bmatrix}
d\theta & d\phi
\end{bmatrix}
\begin{bmatrix}
g_{\theta\theta} & g_{\theta\phi} \\
g_{\phi\theta} & g_{\phi\phi}
\end{bmatrix}
\begin{bmatrix}
d\theta \\
d\phi
\end{bmatrix}.
\]

Therefore the Riemannian metric in local coordinates \((\theta,\phi)\) is
\[
ds^2 = g_{\theta\theta}\, d\theta^2 + 2 g_{\theta\phi}\, d\theta\, d\phi
+ g_{\phi\phi}\, d\phi^2
= d\theta^2 + \sin^2\theta \, d\phi^2.
\]

\textbf{Geometric meaning:}  
The coefficient of \(d\theta^2\) shows that moving along a meridian 
has true arc length proportional to the change in \(\theta\).  
The coefficient \(\sin^2\theta\) shows that circles of latitude 
have circumference \(2\pi\sin\theta\) on the unit sphere.

Thus, \((S^2, g)\) is a Riemannian manifold: the smooth surface of the sphere together 
with the induced inner product structure that varies smoothly from point to point.
\end{example}



\begin{definition}[Submanifold and embedded submanifold]
A subset $N$ of $M$ is a submanifold of $M$ if it is included in $M$ and is a manifold. The submanifold $N$ of $M$ is
called an embedded submanifold of $M$ if there exists a smooth manifold $X$ and a smooth embedding f such that:
$N = f(X)$.
\end{definition}

The objective of the rVAE is submanifold learning.

The Riemannian metric $g$ of $M$ induces an infinitesimal volume element on each tangent space. Thus, a measure on
the manifold $M$ has the expression $dM = \sqrt{\textrm{det}(g(x))}dx$ in local coordinate system $x$.

\begin{example}[Volume element on the 2-sphere]
On the unit sphere in spherical coordinates $(\theta,\phi)$ the metric is
\[
ds^2 = d\theta^2 + \sin^2\theta\, d\phi^2,
\]
so the metric matrix is
\[
g(\theta,\phi) =
\begin{bmatrix}
1 & 0 \\
0 & \sin^2\theta
\end{bmatrix}.
\]
Its determinant is
\[
\det(g) = \sin^2\theta,
\]
hence the Riemannian volume element is
\[
dM = \sqrt{\det(g)}\, d\theta\, d\phi = \sin\theta\, d\theta\, d\phi.
\]
Integrating over $\theta\in[0,\pi],\ \phi\in[0,2\pi]$ yields
\[
\int_0^{2\pi}\int_0^\pi \sin\theta\, d\theta\, d\phi = 4\pi,
\]
the total surface area of the unit sphere.
\end{example}


\begin{definition}[Weighted Submanifold] 
Given a complete $d$-dimensional Riemannian manifold $(M, g)$ and a smooth probability distribution $\omega: M \rightarrow \mathbb{R}$, the weighted manifold $(M, \omega)$ associated to $M$ and $\omega$ is defined as the triplet:

\begin{equation}
   (M, g, \omega \cdot dM) 
\end{equation}

where $dM$ denotes the Riemannian measure of $M$.
\end{definition}


\begin{example}[Weighted sphere]
Consider the unit 2-sphere
\[
S^2 = \{ (x,y,z) \in \mathbb{R}^3 : x^2+y^2+z^2 = 1 \},
\]
equipped with the standard Riemannian metric induced from $\mathbb{R}^3$.  
The corresponding Riemannian volume element in spherical coordinates $(\theta,\phi)$ is
\[
dM = \sin\theta \, d\theta\, d\phi.
\]

Now define a smooth probability density $\omega(\theta,\phi) = \tfrac{1}{4\pi}$, the uniform distribution on the sphere.  
Then the weighted manifold is given by
\[
\bigl(S^2, g, \omega \cdot dM \bigr)
= \bigl(S^2, g, \tfrac{1}{4\pi} \sin\theta \, d\theta \, d\phi \bigr).
\]

In this case, $\omega \cdot dM$ is exactly the uniform probability measure on the sphere.
\end{example}

\subsection{Geodesic and Geodesic Subspaces}
Let $u, v$ be vector fields on $M$. The Riemannian metric $g$ on $M$ induces the notion of covariant derivative. 
The covariant derivative of $v$ in the direction of $u$, written $\nabla_u v$, represents the change of the vector field $v$
in the $u$ direction with the normal component subtracted. In other words, it is the orthogonal projection of the Euclidean directional 
derivative onto the manifolds tangent space.


But it doesn’t mean the result necessarily points in the direction of $u$. Think of $\nabla_u v$ answering: “How does the field 
$v$ change when I walk a little step in direction $u$?” The outcome is another tangent vector, but its orientation can be arbitrary in the tangent space. For instance, on the 2-sphere, parallel transporting a vector along a latitude line (a small circle) produces a derivative that points in a different tangent direction than $u$.


Let $\gamma$ be a curve on $M$ parameterized by $t$, as $\gamma:[0,1] \mapsto M, t \mapsto \gamma(t)$. Its velocity is: $\dot{\gamma}=\frac{d \gamma}{d t}$. Let $u$ be a vector field $t \mapsto u(t)$ defined along $\gamma$. For a vector field $u(t)$ defined along $\gamma$, the covariant derivative of $u$ in the direction of velocity $\dot{\gamma}$ is $\nabla_{\dot{\gamma}} u$.

\begin{definition}[Geodesic]
The curve $\gamma:[0,1]\to M$ is called a geodesic if it satisfies
the equation
\[
\nabla_{\dot{\gamma}} \dot{\gamma} = 0,
\]
where $\nabla$ denotes the covariant derivative of the Levi--Civita
connection associated with the Riemannian metric of $M$.
\end{definition}

A geodesic on $S^2$ is, by definition, a curve  $\gamma(t)$ on the sphere such that its velocity vector 
$\dot{\gamma}(t)$ is transported parallel to itself:
\[
\nabla_{\dot{\gamma}} \dot{\gamma} = 0.
\]
Intuitively, this means the curve has no extrinsic acceleration tangent to the sphere---it ``moves straight'' within the curved surface. 
In other words, this means “move forward without turning”--it’s the generalization of a straight line to curved spaces. 

A geodesic does not bend intrinsically on the surface: its velocity  vector is parallel transported along itself, making it the ``straightest'' 
possible curve with respect to the surface geometry. However, because  the surface itself is curved, the geodesic may appear bent when viewed 
extrinsically in the ambient Euclidean space.


\begin{example}
Consider the unit sphere $S^2 \subset \mathbb{R}^3$ equipped with the Riemannian metric induced from the standard Euclidean metric on $\mathbb{R}^3$.
Let $\gamma(t) = (\cos t, \sin t, 0)$ be a parametrized curve on the sphere, which traces the equator.

The velocity of $\gamma$ is
\[
\dot{\gamma}(t) = (-\sin t, \cos t, 0).
\]
To compute the covariant derivative $\nabla_{\dot{\gamma}} \dot{\gamma}$, we first look at the Euclidean derivative of $\dot{\gamma}$:
\[
\frac{d}{dt}\dot{\gamma}(t) = (-\cos t, -\sin t, 0).
\]

This derivative is not tangent to the sphere: in fact,
\[
\frac{d}{dt}\dot{\gamma}(t) = -\gamma(t),
\]
which points points radially inward, orthogonal to the tangent space of the sphere at $\gamma(t)$. The covariant derivative $\nabla_{\dot{\gamma}} \dot{\gamma}$ is obtained by projecting the Euclidean derivative back onto the tangent space of $S^2$.
Since the derivative is purely normal to $S^2$, its tangential component vanishes:
\[
\nabla_{\dot{\gamma}} \dot{\gamma} = 0.
\]

Hence the equator is a geodesic. More generally, any great circle (i.e.\ the intersection of $S^2$ with a plane through the origin)
satisfies $\nabla_{\dot{\gamma}} \dot{\gamma}=0$ and is therefore a geodesic. Small circles, such as the parallels of latitude that are not great circles, do not satisfy this condition and are not geodesics.
\end{example}

Intuitively, a geodesic is a curve that is parallel to itself. In this sense, a geodesic generalizes the notion of linear
curve on Euclidean spaces. As an example, geodesics on the sphere are the great circles. Geodesics are the equivalent on Riemannian manifolds
of 1D linear subspaces in Euclidean spaces. We define “geodesic subspaces”, that generalize to manifolds higher dimensional
subspaces of Euclidean spaces.

\paragraph{Christoffel Symbols} 
After differentiating a basis vector, projecting onto the tangent plane, and decomposing into basis vector components, we obtain the Christoffel symbols. They are written as

\[
\nabla_{\partial_i} \vec{e}_j 
= \sum_{k \in \{\theta,\phi\}} \Gamma^k_{ij}\,\vec{e}_k.
\]

Here:
\begin{itemize}
  \item $k$ indicates the basis vector direction the result points along after projection (the component we are looking at),
  \item $i$ indicates the direction of differentiation (the variable with respect to which we differentiate),
  \item $j$ indicates the basis vector being differentiated.
\end{itemize}

\begin{example}
Consider the basis vector
\[
\vec{e}_\theta = \frac{\partial \vec{X}}{\partial \theta} 
= (\cos\theta \cos\phi,\ \cos\theta \sin\phi,\ -\sin\theta).
\]

Differentiate with respect to $\phi$:
\[
\frac{\partial \vec{e}_\theta}{\partial \phi} 
= (-\cos\theta \sin\phi,\ \cos\theta \cos\phi,\ 0).
\]


Now project this derivative onto the tangent plane spanned by 
$\{\vec{e}_\theta, \vec{e}_\phi\}$.

Dot product with $\vec{e}_\theta$ gives $0$.
Dot product with $\vec{e}_\phi = (-\sin\theta \sin\phi,\ \sin\theta \cos\phi,\ 0)$ gives $\cos\theta \sin\theta$.

Since $\|\vec{e}_\phi\|^2 = \sin^2\theta$, the coefficient is
\[
\frac{\cos\theta \sin\theta}{\sin^2\theta} = \cot\theta.
\]

So the tangential decomposition is
\[
\nabla_{\partial_\phi}\vec{e}_\theta = \cot\theta\, \vec{e}_\phi.
\]

Written in Christoffel symbol form,
\[
\nabla_{\partial_\phi}\vec{e}_\theta 
= \Gamma^\theta_{\phi\theta}\,\vec{e}_\theta 
+ \Gamma^\phi_{\phi\theta}\,\vec{e}_\phi,
\]
with
\[
\Gamma^\theta_{\phi\theta} = 0,
\qquad
\Gamma^\phi_{\phi\theta} = \cot\theta.
\]

\end{example}


\paragraph{Christoffel Symbols from the Metric}

We begin from the definition of the metric:
\[
g_{ij} = \vec e_i \cdot \vec e_j.
\]

Differentiating with respect to $x^k$ gives
\[
\partial_k g_{ij}
= (\partial_k \vec e_i)\cdot \vec e_j
+ \vec e_i \cdot (\partial_k \vec e_j).
\]

Using the covariant derivative expansion
\[
\nabla_{\partial_k}\vec e_i
= \Gamma^p_{ki}\, \vec e_p,
\]

we can replace $\partial_k \vec e_i$ by
\[
\partial_k \vec e_i = \Gamma^p_{ki}\, \vec e_p.
\]

Thus
\[
\begin{aligned}
\partial_k g_{ij} &= \Gamma^p_{ki}\,\vec e_p\!\cdot\!\vec e_j \;+\; \vec e_i\!\cdot\!\Gamma^p_{kj}\,\vec e_p\\
&= \Gamma^p_{ki}\, g_{pj} \;+\; \Gamma^p_{kj}\, g_{ip}
\end{aligned}
\]


Now, we write down the three cyclic permutations:

\[
\partial_k g_{lm} = \Gamma^p_{kl} g_{pm} + \Gamma^p_{km} g_{lp},
\]
\[
\partial_l g_{km} = \Gamma^p_{lk} g_{pm} + \Gamma^p_{lm} g_{kp},
\]
\[
\partial_m g_{kl} = \Gamma^p_{mk} g_{pl} + \Gamma^p_{ml} g_{kp}.
\]

Adding the first two, subtracting the third and using $\Gamma^p_{mk} = \Gamma^p_{km}$, $\Gamma^p_{ml} = \Gamma^p_{lm}$ 
$\Gamma^p_{kl} = \Gamma^p_{lk}$ (Levi–Civita connection is torsion-free), $g_{lp} = g_{pl}$ (metric tensor is symmetric) we get:
\[
\partial_k g_{lm} + \partial_l g_{km} - \partial_m g_{kl}
= 2\,\Gamma^p_{kl}\, g_{pm}.
\]

Finally, contracting both sides with the inverse metric $g^{im}$ yields
\[
\Gamma^i_{kl}
= \tfrac{1}{2} g^{im}\big(
\partial_k g_{lm}
+ \partial_l g_{km}
- \partial_m g_{kl}
\big).
\]

This is the standard formula for the Christoffel symbols of the Levi--Civita connection.


\begin{definition}[Parallel transport]
Let $\gamma:[0,1]\to M$ be a smooth curve on a Riemannian manifold $(M,g)$. A vector field $u(t)$ along $\gamma$ is called \emph{parallel transported}
if it satisfies
\[
\nabla_{\dot\gamma(t)} u(t) = 0.
\]
This means the vector field does not twist relative to the manifold's geometry as it moves along the curve.
\end{definition}


\begin{example}[Parallel transport on the sphere]
On the unit sphere $S^2 \subset \mathbb{R}^3$, consider a tangent vector at a point on the equator. If the curve $\gamma$ is the equator itself 
(a geodesic), then parallel transport keeps the vector pointing in the  same direction relative to the equator.
\end{example}

\subsection{Riemannian Exp and Log Maps}

For any point $x \in M$ and tangent vector at this point $u \in T_xM$, there exits a unique geodesic $\gamma$, with initial conditions 
$\gamma(0) = x$ and $\gamma'(0) = u$. In general, the existence of the geodesic is only guaranteed locally around 0. As we assume 
the completeness of the manifold $M$, the geodesics are defined on $\mathbb{R}$.

\begin{definition}[Riemannian exp map]
Let $\gamma$ be the unique geodesic with initial condition $\gamma(0)=x$ and  $\gamma'(0)=u$. We define the Riemannian exponential map 
at $x$ as:

\begin{equation}
    \mathrm{Exp}_x(u) = \gamma(1)
\end{equation}
\end{definition}


The exponential map takes a point $x$ and tangent vector $u$ and returns the point at time 1 obtained by “shooting” with
$u$ along the geodesic. The exponential map is a local diffeomorphism onto a neighbourhood of $x$ in $M$. Let $V(x)$
be the largest neighbourhood on which the exponential map at $x$ is a diffeomorphism. We can define its inverse on $V (x)$.


\begin{definition}[Riemannian log map]
The Riemannian exponential map has an inverse on the domain $V (x)$, defined as the Riemannian log map at $x$, and denoted $Log_x$.
\end{definition}

The domain $V (x)$ is the global bijectivity domain of the exponential map at $x$. 
One can show that $V (x)$ is a star-shaped domain delimited by a continuous curve $\mathcal{C}_x$.

\subsection{Tangential cut locus, cut locus and Injectivity radius}

\begin{definition}
    The curve $\mathcal{C}_x$ delimiting $V(x)$ is called the tangential cut locus of $x$. 
    The image of $\mathcal{C}_x$ by the exponential map at $x$ is called the cut locus of $x$.
\end{definition}

We can define the distance from a point $x$ to its cut locus $C = \mathcal{C}_x$.

In 2-sphere, $\mathcal{C}_x$ is the circle of radius $\pi$ in the tangent plane (tangential cut locus).
Its image under $\textrm{Exp}_x$ is the antipodal point (the cut locus).


On the unit 2-sphere $S^2$, the exponential map $\mathrm{Exp}_x$ at a point $x \in S^2$ takes a tangent vector $u \in T_xS^2$ 
and maps it to the point reached by following the geodesic in the direction of $u$ for length $\|u\|$. 

The injectivity radius at $x$ is the largest radius $r$ such that $\mathrm{Exp}_x$ is one-to-one on the ball $B_r(0) \subset T_xS^2$. 
On $S^2$, geodesics starting from $x$ remain unique until they meet again at the antipodal point, a distance $\pi$ away. 
Thus
\[
\mathrm{inj}_x(S^2) = \pi \quad \text{for all } x \in S^2.
\]

Since this value is the same for every point, the global injectivity radius is
\[
\mathrm{inj}(S^2) = \pi.
\]

Notice that the diameter of $S^2$ is also $\pi$, so here we achieve equality in the general bound
\[
0 < \mathrm{inj}(M) \leq \mathrm{diam}(M).
\]


\subsection{Hadamard Manifolds}

\begin{theorem}[Cartan--Hadamard]
Let $(M,g)$ be a simply connected, complete Riemannian manifold with non-positive sectional curvature (i.e.\ a Hadamard manifold). 
Then, for any point $x \in M$, the exponential map
\[
\mathrm{Exp}_x : T_xM \;\to\; M
\]
is a global diffeomorphism. 
\end{theorem}

\begin{corollary}
On a Hadamard manifold, the logarithm map
\[
\mathrm{Log}_x : M \;\to\; T_xM
\]
is globally well-defined for every $x \in M$. 
In particular, every point on the manifold corresponds uniquely to a tangent vector at $x$, and computations on $M$ can be transferred to the Euclidean vector space $T_xM$ via the logarithm map.
\end{corollary}

\subsection{Distance in $M$ and between submanifolds of $M$}

\begin{definition}[Riemannian distance]
For any point $y \in V(x)$, the Riemannian distance function is given by
\[
    d_M(x, y) = \| \mathrm{Log}_x(y) \|_{g(x)}
\]
where the norm corresponds to the inner product $g(x)$ at $x$.
\end{definition}

\begin{definition}[Wasserstein distance]
The 2-Wasserstein distance between probability measures $\mu$ and $\nu$ defined on $M$ is given by
\[
    d_2(\mu, \nu) =
    \left(
        \inf_{\gamma \in \Gamma(\mu, \nu)}
        \int_{M \times M}
        d_M^2(x_1, x_2) \, d\gamma(x_1, x_2)
    \right)^{1/2},
\]
where $\Gamma(\mu, \nu)$ denotes the collection of all measures on $M \times M$ with marginals $\mu$ and $\nu$ on the first and second factors, respectively. A transport plan $\gamma$ is a probability measure on $M \times M$ that tells how much mass from point $x_1$ is moved to point $x_2$.
\end{definition}

This notion allows defining a distance between two embedded submanifolds of $M$, with the Riemannian metric associated to their embedding $f$. We consider the 2-Wasserstein distance associated with the corresponding distributions, interpreted as singular distributions in $M$. This notion of distance is the evaluation metric used to compare weighted submanifold learning techniques.


\subsection{Fr\'echet Mean}

\begin{definition}[Fr\'echet Mean]
Let $X$ be a random variable taking values on a manifold $M$ with probability distribution $p$.

The \emph{population Fr\'echet mean} is defined as
\[
\mu = \arg\min_{y \in M} \int_M d_M^2(x,y) \, dp(x).
\]

Given a random sample $X_1,\dots,X_n$ on $M$, the \emph{sample Fr\'echet mean} is defined as
\[
\bar x = \arg\min_{y \in M} \frac{1}{n} \sum_{i=1}^n d_M^2(y, x_i).
\]
\end{definition}

\section{Proof of Proposition 1}

\subsection{Preliminaries}

We begin by introducing notations. For $T \in \mathbb R^*_+$, we denote by $\mu_T$ the standard multivariate normal in $\mathbb{R}^L$ truncated at a distance $T$ from the origin. We write $B_T$ for the closed ball of $\mathbb{R}^L$ of radius $T$, which is a compact subset of $\mathbb{R}^L$ as it is bounded and closed in $\mathbb{R}^L$.

We write $C(\mathbb{R}^L,\mathbb{R}^D)$ for the set of continuous functions from $\mathbb{R}^L$ to $\mathbb{R}^D$, and we write $C(B_T,\mathbb{R}^D)$ for the set of continuous functions from $B_T$ to $\mathbb{R}^D$.

We first introduce a lemma that generalizes a result from~\cite{hornik1991approximation} to functions in $D$ dimensions.

\begin{lemma}\label{lemma:d_approx}
Consider $K$ a compact subset of $\mathbb{R}^L$. For any $\epsilon > 0$ and any continuous function $f \in C(K,\mathbb{R}^D)$, there exists a neural network represented by the function $f_\theta$, parameterized by $\theta$, such that:
\[
\sup_{z \in K} \| f(z) - f_\theta(z) \|_2^2 < D\epsilon^2.
\]
\end{lemma}

\begin{proof}
Take $\epsilon > 0$. For each $d \in \{1,\dots,D\}$, consider $f_d$, the $d$-th component of $f \in C(K,\mathbb{R}^D)$. Then $f_d \in C(K,\mathbb{R})$. From the results of \cite{hornik1991approximation}, there exists a neural network that approximates $f_d$ with precision $\sqrt\epsilon$.

In our notations, $L$ is the dimension of the input space (as opposed to $k$ in \cite{hornik1991approximation}), $z \in \mathbb{R}^L$ is an input element, and $K$ is the number of hidden units (as opposed to $n$ in \cite{hornik1991approximation}). We define:
\[
M^{(K)}_L(\psi) = \Big\{ h : \mathbb{R}^L \to \mathbb{R} \,\big|\,
h(z) = \sum_{k=1}^K \beta_k \, \psi(a_k'z - \theta_k) \Big\},
\]
the space of functions represented by a fully connected neural network with one output unit and one hidden layer with $K$ hidden units, where $\psi$ is the activation function. Then define
\[
M_L(\psi) = \bigcup_{K=1}^{+\infty} M^{(K)}_L(\psi),
\]
the space of functions represented by one-hidden-layer fully connected networks with any number of hidden units. We assume $\psi$ is continuous, bounded, and nonconstant (for example, the sigmoid).

By Theorem 2 of \cite{hornik1991approximation}, $M_L(\psi)$ is dense in $C(K,\mathbb{R})$. Thus, for each $f_d \in C(K,\mathbb{R})$, there exists $h_d \in M_L(\psi)$ such that:
\[
\sup_{z \in K} |f_d(z) - h_d(z)| < \epsilon.
\]

We fix functions $h_1,\dots,h_D$ approximating the $D$ components of $f$. We define $h = (h_1,\dots,h_D) : \mathbb{R}^L \to \mathbb{R}^D$. Then
\begin{align*}
\sup_{z\in K}\,\| f(z) - h(z) \|_2^2
&= \sup_{z\in K}\,\sum_{d=1}^D |f_d(z) - h_d(z)|^2\\
&\le \sum_{d=1}^D \sup_{z\in K} |f_d(z) - h_d(z)|^2\\
&< D\,\epsilon^2.
\end{align*}


Thus, $f$ can be approximated by the neural network $h$ with precision $\epsilon$, obtained by combining the networks defining $h_d$, $d=1,\dots,D$.
\end{proof}

\begin{proposition}

Let \((N_T,\nu_T)\) be a weighted Riemannian submanifold of \(M\), embedded in a submanifold \(L\subset M\) that is homeomorphic to \(\mathbb{R}^L\). Suppose there exists an embedding \(f^M:\mathbb{R}^L\to L\) such that \(\nu_T=(f^M)_*\mu_T\), where \(\mu_T\) is the standard Gaussian on \(\mathbb{R}^L\) truncated to the ball \(B_T\). Assume there exists \(\mu\in M\) with \(N\subset V(\mu)\), where \(V(\mu)\) is the maximal domain on which the Riemannian exponential map \(\mathrm{Exp}_\mu^M\) is a global bijection. Then, for any \(0<\epsilon<1\), there exists a Riemannian VAE with decoder given by a neural network \(f^M_\theta\) (parameterized by \(\theta\)) such that

\[
d_2^2\bigl(N_T,\,N_{\theta,T}\bigr)\;<\;C^2_{T,D}\,D\epsilon^{2},
\]
where \(d_2\) denotes the \(2\)-Wasserstein distance, \(C_{T,D}>0\) depends only on \(T\) and the dimension \(D\), and

\[
    (N_{\theta,T}, \nu_{\theta,T}) \;=\;\bigl(f^M_\theta(\mathbb{R}^L),\, (f^M_\theta)_* \mu_T\bigr),
\]

\end{proposition}

\paragraph{Explanation}
We have a target geometric--probabilistic object \((N_T,\nu_T)\): a submanifold \(N_T\subset M\) together with a probability measure \(\nu_T\) supported on it (a “weighted submanifold”). It sits inside another submanifold \(L\subset M\) that is homeomorphic (the same shape topologically) to \(\mathbb{R}^L\). This supplies a convenient flat parameter space. There is an embedding \(f^M:\mathbb{R}^L\to L\) whose pushforward of a truncated standard Gaussian \(\mu_T\) equals \(\nu_T\):
\[
(f^M)_*\mu_T=\nu_T.
\]
Equivalently, if \(z\sim \mu_T\) in \(\mathbb{R}^L\) and \(x=f^M(z)\), then \(x\) has distribution \(\nu_T\).

We assume there exists a base point \(\mu\in M\) such that \(N\subset V(\mu)\), where \(V(\mu)\) is the maximal domain on which the Riemannian exponential map \(\mathrm{Exp}_\mu^M\) is a global bijection. Practically, geodesic coordinates from \(\mu\) are well behaved on the whole set (no self-overlap, no cut locus).


For any accuracy level \(0<\epsilon<1\), there exists a Riemannian VAE with decoder given by a neural network \(f^M_\theta:\mathbb{R}^L\to L\) 
with $L \subset M$  such that the model distribution

\[
    (N_{\theta,T}, \nu_{\theta,T}) \;=\;\bigl(f^M_\theta(\mathbb{R}^L),\, (f^M_\theta)_* \mu_T\bigr),
\]


is close to the target \((N_T,\nu_T)\) in quadratic Wasserstein distance:
\[
d_2^2\!\bigl(N_T,\,N_{\theta,T}\bigr)\;<\; C_{T,D}^2\,D\,\epsilon^{2}.
\]
Here \(C_{T,D}>0\) depends only on the truncation radius \(T\) and the output dimension \(D\).

The two roles of \(f^M_\theta\) are:

\begin{enumerate}
  \item \textbf{Geometric mapping:} 
  \[
    f^M_\theta:\mathbb{R}^L \;\longrightarrow\; L \subset M.
  \]
  Its image is a submanifold inside \(M\), the candidate learned submanifold. 
  In the notation
  \[
    (N_{\theta,T}, \nu_{\theta,T}) \;=\;\bigl(f^M_\theta(\mathbb{R}^L),\, (f^M_\theta)_* \mu_T\bigr),
  \]
  the first component is the geometric support.

  \item \textbf{Measure transport:} 
  At the same time, \(f^M_\theta\) pushes forward the latent distribution 
  \(\mu_T\) (the truncated Gaussian in \(\mathbb{R}^L\)) to a distribution on \(L\):
  \[
    (f^M_\theta)_* \mu_T.
  \]
  This is the probabilistic weighting on the learned manifold.
\end{enumerate}



The embedding \(f^M\) already maps the simple latent \(\mu_T\) to the exact target \(\nu_T\). By universal approximation, a neural decoder \(f^M_\theta\) can approximate \(f^M\) uniformly well on the compact ball \(B_T\subset\mathbb{R}^L\). Uniform control on the map, combined with the tame geometry ensured by the exponential-chart assumption, yields control of transport cost, hence the  Wasserstein distance bound with a constant–times–squared–error form.


\begin{proof}
We can write $L = f^M(\mathbb{R}^L)$ the embedded submanifold of $M$, that contains $N_T: N_T \subset L = f^M(\mathbb{R}^L)$.
Without loss of generality, we can restrict $N_T$ to the support of the singular probability distribution $\nu_T$.
We write $N_T = f^M(B_T)$. As $B_T$ is compact and $f^M$ is continuous by definition of an embedding, we observe that
$N_T$ is compact.

As $N_T\subset V(\mu)$, we can define the function:

\[
\begin{aligned}
f\colon B_T &\mapsto T_\mu M\\
z &\mapsto f(z) = \mathrm{Log}_\mu f^M(z).
\end{aligned}
\]


So, $f^M$ maps $\mathbb{R}^L$ to $L$ and $f$ maps it to tangent space at $\mu$ point of $M$.

We have $f \in C(B_T,\mathbb{R}^D)$, where $B_T$ is a compact subset of $\mathbb{R}^L$. From Lemma~\ref{lemma:d_approx}, there exists 
a neural network represented by $f_\theta$, such that:

\[
\sup_{z \in B_T} \| f(z) - f_\theta(z) \|_2^2 < D\epsilon^2.
\]

We introduce the weighted submanifold $(N_{\theta,T},\nu_{\theta,T})$ associated to the neural network 
represented by $f_\theta$:

\[
(N_{\theta,T},\nu_{\theta,T}) \;=\;\bigl(f^M_\theta(B_T),\, (f^M_{\theta}|_{B_T})_* \mu_T\bigr),
\]

which corresponds to the generative model defining the Riemannian VAE. We show that $N_{\theta,T}$ can approximate $N_T$
with a precision $\epsilon$, in terms of the 2-Wasserstein distance.

By definition, the squared 2-Wasserstein distance between weighted submanifolds is the 2-Wasserstein distance between 
the probability measures defined on them. Here, the probability distributions are $\nu_T$ and $\nu_{\theta,T}$, so we
write:

\[
    d_2^2(N_T,N_{\theta,T}) =
        \inf_{\gamma \in \Gamma(\nu_T, \nu_{\theta,T})}
        \int_{N_T \times N_{\theta,T}}
        d_M^2(x_1, x_2) \, d\gamma(x_1, x_2).
\]

Now, we consider the homeomorphism:

\[
\bigl((f^M)^{-1}, (f^M_\theta)^{-1}\bigr) : N_T \times N_{\theta,T} \mapsto B_T \times B_T
\]

which is a “global chart” of the submanifold $N_T \times N_{\theta,T}$ of $M \times M$.

Now, we perform the change of variables $(x_1, x_2)=(f^M(z_1), f^M_\theta(z_2))$ to represent 
point $(x_1,x_2)$. We get:

\[
    d_2^2(N_T,N_{\theta,T}) =
        \inf_{\gamma \in \Gamma(\mu_T, \mu_T)}
        \int_{B_T \times B_T}
        d_M^2(f^M(z_1), f^M_\theta(z_2)) \, d\gamma(z_1, z_2),
\]

where the probability density $\gamma$ is expressed in this chart:

\[
 d\gamma(z_1, z_2) = \gamma(z_1, z_2)dz_1dz_2,
\]

and we have $((f^M_\theta)^{-1})_*\nu_{\theta,T} = \mu_T$. We consider the following element $\tilde{\gamma}$ of $\Gamma(\mu_T, \mu_T)$:

\[
 \tilde{\gamma}(z_1,z_2) = \mu_T(z_1)\delta_{z1=z2} = \mu_T(z_2)\delta_{z1=z2}.
\]

So, $\tilde{\gamma}$ is the identity coupling: all the mass at $z_1$ in the first copy is transported to exactly the same point 
$z_2$ in the second copy and it is the canonical “do nothing” transport plan.

By the property of the $\mathrm{inf}$, we get the upper bound:

\[
\begin{aligned}
d_2^2(N_T,N_{\theta,T})
   &\le \int_{B_T \times B_T}
        d_M^2\!\bigl(f^M(z_1), f^M_\theta(z_2)\bigr) \, d\tilde{\gamma}(z_1, z_2) \\[6pt]
   &= \int_{B_T}
        d_M^2\!\bigl(f^M(z), f^M_\theta(z)\bigr) \, d\mu_T(z).
\end{aligned}
\]

We define compact $K$ of $T_\mu M$ as the compact $f(B_T)$ extended in $L_2$ norm by a distance $\sqrt D$ in 
all directions. We have, for any $z \in B_T : f(z) \in f(B_T) \subset K$ and $f_\theta(z) \in K$. 
The Riemannian exponential $\mathrm{Exp}(\mu,\bullet)$ is continuous on the compact $K$, and therefore Lipshitz on $K$. 
There exists constant $C_{T,D}$ such that for any $y_1, y_2 \in K$, we have:

\[
d_M(\mathrm{Exp}(\mu,y_1), \mathrm{Exp}(\mu,y_2)) < C_{T,D}\|y_1-y_2\|_2.
\]

In simple terms: if we start at the same base point $\mu$, take two tangent vectors 
$y_1, y_2$, and map them to the manifold with the exponential map, then the distance 
between the two resulting points on the manifold cannot grow faster than a constant 
multiple of the Euclidean distance between the vectors in the tangent space.

So the exponential map behaves in a ``well-controlled'' way on this compact region: 
it will not blow up small differences in tangent space into huge differences on the 
manifold. The constant $C_{T,D}$ is the quantitative bound that guarantees this control.

We apply this to $y_1 = f(z)$ and $y_2 = f_\theta(z)$ for any $z \in B_T$, and taking square:

\[
\begin{aligned}
d_M^2(\mathrm{Exp}(\mu,f(z)), \mathrm{Exp}(\mu,f_\theta(z))) &< C^2_{T,D}\|f(z)- f_\theta(z)\|^2_2\\[6pt]
\implies d_M^2(f^M(z), f^M_{\theta}(z)) &< C^2_{T,D}\|f(z)- f_\theta(z)\|^2_2\\[6pt]
\implies d_M^2(f^M(z), f^M_{\theta}(z)) &< C^2_{T,D}\sup_{z \in B_T} \| f(z) - f_\theta(z) \|_2^2\\[6pt]
\implies d_M^2(f^M(z), f^M_{\theta}(z)) &< C^2_{T,D}D\epsilon^2.
\end{aligned}
\]

We integrate left side using the measure $\mu_T$, and write:

\[
\begin{aligned}
d_2^2(N_T,N_{\theta,T}) &< \int_{\mathbb{R}_L} d_M^2(f^M(z), f^M_{\theta}(z))d\mu_T(z)\\[6pt]
&< \int_{\mathbb{R}_L} C^2_{T,D}D\epsilon^2 d\mu_T(z)\\[6pt]
&< C^2_{T,D}D\epsilon^2.
\end{aligned}
\]

Therefore, for any $T$, any weighted submanifold $N_T$ verifying the assumptions can be approximated by a submanifold
$N_{\theta,T}$ generated by the model of a Riemannian VAE.

\end{proof}
\bibliographystyle{unsrt}
\bibliography{references}

\end{document}